% ------------------------------------------------------------
% A7.0b Projection Domain & Stability Conditions (v1.0)
% Author: Borealis S. Hedling
% Altitude: A7.0 (Non-Dual Layer)
% Purpose: Define the domain, admissibility, locality, stability,
%          kernel interaction, and metadata conditions required
%          for safe non-dual projection.
% ------------------------------------------------------------

\documentclass[11pt]{article}
\usepackage{amsmath, amssymb, amsthm}

\title{A7.0b Projection Domain and Stability Conditions}
\author{Borealis S. Hedling}
\date{October 2026}

\begin{document}
\maketitle

\section*{Overview}

This document defines the structural conditions required for safe execution
of A7.0 non-dual projection operators. These conditions specify the domain
of projection, admissibility constraints, locality bounds, stability envelopes,
kernel interaction rules, and metadata hooks required for A7.1 classification
and A7.2 repair.

% ------------------------------------------------------------
\section{Projection Domain}

Non-dual projection is defined only on the A6 manifold:
\[
\mathrm{Dom}(\Pi_{\mathrm{ND}}) = M_{\mathrm{A6}}.
\]

Projection is prohibited on:
\begin{itemize}
  \item regions with undefined curvature,
  \item regions with unbounded holonomy,
  \item regions with spectral indeterminacy,
  \item regions lacking modal truth values.
\end{itemize}

Admissible domain:
\[
M_{\mathrm{adm}} = \{ p \in M_{\mathrm{A6}} : \Gamma(p) \text{ finite} \}.
\]

% ------------------------------------------------------------
\section{Projection Admissibility Conditions}

Projection is admissible only when:
\[
|K(p)| < K_{\max},
\]
\[
\delta_{\mathrm{spec}}(p) < \delta_{\max},
\]
\[
H(\gamma) < \Theta_{\max},
\]
\[
\mu(p) \text{ defined}.
\]

These conditions ensure projection does not fracture the manifold.

% ------------------------------------------------------------
\section{Projection Locality Constraints}

Define the locality radius:
\[
r_{\mathrm{loc}}(p) = \inf \{ r : B_r(p) \cap \Sigma \neq \emptyset \}.
\]

Projection must satisfy:
\[
\Pi_{\mathrm{ND}}(q) = \Pi_{\mathrm{ND}}(p)
\quad \text{for } q \in B_{r_{\mathrm{loc}}}(p).
\]

This prevents non-local collapse propagation.

% ------------------------------------------------------------
\section{Projection Stability Envelope}

Define the stability envelope:
\[
E_{\mathrm{proj}}(p) = E_{\mathrm{A6}}(p) - \Delta E(p),
\]
where:
\[
\Delta E(p) \ge 0.
\]

Projection is stable if:
\[
E_{\mathrm{proj}}(p) > E_{\min}.
\]

This ensures projection does not destabilize the manifold.

% ------------------------------------------------------------
\section{Projection Kernel Interaction}

Non-dual projection breaks kernel invariants:
\[
K_{\mathrm{A6}}(p) \not\Rightarrow K_{\mathrm{A7}}(\Pi_{\mathrm{ND}}(p)).
\]

However, projection must record kernel interaction metadata:
\[
\kappa(p) = K_{\mathrm{A6}}(p).
\]

This metadata is required for A7.2 repair.

% ------------------------------------------------------------
\section{Projection Metadata Hooks}

Projection must emit the following metadata for A7.1 classification:
\[
\Gamma(p) = (K(p), \mathrm{Spec}(p), \mu(p), H(p)).
\]

Additionally, projection must record:
\begin{itemize}
  \item projection locality radius $r_{\mathrm{loc}}(p)$,
  \item projection stability envelope $E_{\mathrm{proj}}(p)$,
  \item kernel interaction $\kappa(p)$,
  \item collapse type candidates $\mathcal{T}_{\mathrm{cand}}(p)$.
\end{itemize}

This metadata is required for collapse-surface classification.

% ------------------------------------------------------------
\section{Conclusion}

A7.0b defines the domain, admissibility, locality, stability, kernel
interaction, and metadata conditions required for safe non-dual projection.
These conditions complete the A7.0 altitude and prepare the system for
A7.1 collapse-surface classification and A7.2 repair.

\end{document}
